
\textbf{Synthetic ModelZoo}: We generate 1000 neural network models distributed across three architecture families: 400 CNNs (40\%), 400 Transformers (40\%), and 200 RNNs (20\%). Each model consists of randomly initialized weights with dimensionality 1000, simplified from realistic dimensions to enable rapid prototyping. The dataset is split 70\% train (700 models: CNN+Transformer), 15\% validation (150 models: CNN+Transformer), 15\% test (150 models: 50 CNN, 50 Transformer, 50 RNN). RNN models appear only in the test set to evaluate frozen-K generalization.

We use synthetic random weights as a proof-of-concept simplification. If the approach cannot succeed on synthetic data with clear structure, it will not succeed on real models with additional complexity. The clear failure signal across all metrics validates this choice—the failure is not due to insufficient realism but mechanism inadequacy.

\textbf{Implementation}: PyTorch 2.0 with single-GPU training. Model architecture: Deep Sets with per-element MLP (input: 1000-dim → hidden: 256-dim → quotient: K=32), architecture embeddings (64-dimensional for {CNN, Transformer, RNN}), mean pooling aggregation, decoder MLP (quotient: K=32 → hidden: 256-dim → output: 1000-dim). Training hyperparameters: Adam (lr=1e-3, weight_decay=1e-4), CosineAnnealingLR (T_max=100 epochs), batch size 32, epochs 20 (early stopped at epoch 12), early stopping patience 10, gradient clipping max_norm=1.0, equivariance loss weight λ_equiv=0.5. Training time: approximately 2 hours for 20 epochs. All experiments use fixed random seeds (seed=42).
